\section{Experiments}
\label{sec:experiments}

\subsection{Experimental setup}

\paragraph{Data and evaluation.}
We use the speech-to-speech retrieval setting introduced by
SpeechDPR~\citep{Lin2024SpeechDPRES}. Spoken questions come from SLUE-SQA-5
and candidate passages are the linked Spoken Wikipedia recordings
\citep{Shon2022SLUEPA}. The official train, validation, test, and
verified-test splits contain 46,186, 1,939, 2,382, and 408 questions. The
current SpeechGR metadata export contains 15,883 unique linked passages and
100\% coverage of the test gold identifiers. Each query has one gold passage.
We report Hit@1 and Hit@20, the percentage of queries for which that passage
appears in the first one or twenty generated identifiers. The headline result
uses the reported test split. The ablations are presented as development-stage
experiments, but the evaluation split and shared pretraining setting are not
recorded for every row. We therefore interpret them as comparisons among the
stated configurations rather than repeated-run or test-set effects.

The audio is natural, but the retrieval semantics are text-derived. Questions
come from five text QA datasets and were read by crowd workers; candidate
passages were paired using transcripts, answer-string matching, BM25,
sentence-level text retrieval, and text QA models
\citep{Shon2022SLUEPA}. Thus, relevance is answer-based rather than acoustic,
but the construction may reward lexical overlap. This makes clean-text
retrieval an important reference while preventing strong claims about
speech-specific semantics or compositional reasoning.

\paragraph{SpeechGR configuration.}
The headline system uses Flan-T5-base, HuBERT-large layer 22, a \(K=500\)
codebook, repetition collapse without BPE, ten epochs of span-corruption
adaptation on 6k hours of LibriLight, cross-entropy training, and beam-20
trie-constrained decoding
\citep{10.5555/3722577.3722647,Hsu2021HuBERTSS,9052942}.
The thesis conflicts on whether the \(k\)-means codebook was fitted on 100 or
960 hours of LibriSpeech; we cannot resolve that provenance from the PDF.

\paragraph{Comparisons.}
We include a ground-truth-text BM25 reference on the same 15,883-passage
metadata export used by SpeechGR. We also report SpeechDPR, BGE, CLSR, and
WavRAG results from prior work rather than reproducing them
\citep{Lin2024SpeechDPRES,hu2026clsr,chen-etal-2025-wavrag}. Their supervision,
retrieval mechanism, and candidate sets differ. SpeechGR is index-free in the
generative-retrieval sense: it uses no external document-embedding index or
approximate-nearest-neighbor search.

\begin{table*}[t]
\centering
\small
\setlength{\tabcolsep}{4pt}
\begin{tabularx}{\textwidth}{@{}l l >{\raggedright\arraybackslash}X l rrr@{}}
\toprule
Model & Query & Text-side signal & Retrieval / candidates & @1 & @10 & @20 \\
\midrule
\multicolumn{7}{@{}l}{\textit{Current 15,883-passage collection}} \\
BM25$^\ast$ & gold text & query and passage transcripts & sparse / 15,883 & \textbf{12.72} & -- & \textbf{43.20} \\
SpeechGR & speech & text-pretrained LM; no task text inputs & GR / 15,883 & 5.81 & -- & 18.452 \\
\midrule
\multicolumn{7}{@{}l}{\textit{Prior reported context; protocols are not matched}} \\
SpeechDPR w/o KD$^\dagger$ & speech & no teacher KD & dense / \(\sim\)39k & -- & -- & 0.04 \\
SpeechDPR with KD$^\dagger$ & speech & UASR + text-retriever KD & dense / \(\sim\)39k & -- & -- & 19.73 \\
BGE clean text$^\dagger$ & gold text & query and passage text & dense / n.r. & 38.71 & 84.34 & -- \\
CLSR$^\dagger$ & speech & ASR losses + frozen BGE space & dense / n.r. & 30.65 & 74.43 & -- \\
WavRAG$^\dagger$ & speech & mixed-modal training incl.\ text/TTS & dense / n.r. & 33.92 & 72.21 & -- \\
\bottomrule
\end{tabularx}
\caption{Test-set retrieval context in percent. With one labelled passage,
Hit@\(k\) and Recall@\(k\) use the same success criterion, but the protocols
and candidate pools differ. \(\ast\) Gold-text evaluation using
\texttt{BM25Okapi} v0.2.2 (\(k_1=1.5,b=0.75,\epsilon=0.25\)) with lowercase
alphanumeric tokenization. \(\dagger\) Values imported from prior work
\citep{Lin2024SpeechDPRES,hu2026clsr,chen-etal-2025-wavrag}, not reproduced.
``n.r.'' denotes a candidate-pool size not reported clearly enough for a
matched comparison.}
\label{tab:main-results}
\end{table*}

\subsection{Main results}

SpeechGR obtains 5.81 Hit@1 and 18.452 Hit@20
(Table~\ref{tab:main-results}). On the same 15,883 passages, the gold-text BM25
reference reaches 12.72 Hit@1 and 43.20 Hit@20. It is not speech-input matched,
but shows that the task strongly rewards lexical access. SpeechGR's non-zero
result shows that discrete units retain enough information to learn some
question-to-passage association; the text-derived evaluation cannot determine
whether that association reflects compositional semantics or lexical matching.

SpeechDPR's 19.73 Hit@20 is also not directly comparable: it uses
text-retriever distillation, external dense passage search, and roughly 39k candidates
rather than 15,883. Its 0.04 Hit@20 without distillation nevertheless isolates
a striking dependence on text-retriever guidance in that system. SpeechGR's
result should therefore be read as feasibility, not near-parity or a
state-of-the-art claim.

Recent systems widen the numerical context. CLSR reports 30.65 R@1 and 74.43
R@10 after transcript-supervised ASR training and mapping speech into a frozen
BGE representation space; WavRAG reports 33.92/72.21 after large-scale
mixed-modality training. Because their candidate pools are under-specified,
we do not compute gaps to SpeechGR. Their results are nevertheless consistent
with the value of text-rich semantic supervision on this text-derived task.

\begin{table*}[t]
\centering
\small
\setlength{\tabcolsep}{5pt}
\begin{tabular}{@{}cllcc@{}}
\toprule
Panel & HuBERT / pretraining setting & Compression & Hit@1 & Hit@20 \\
\midrule
\multicolumn{5}{@{}l}{\textit{A. Acoustic-unit design}} \\
& layer 7, \(K=500\) & none & 1.20 & 6.00 \\
& layer 22, \(K=128\) & none & \textbf{3.76} & 13.04 \\
& layer 22, \(K=500\) & none & 3.61 & \textbf{14.189} \\
& layer 22, \(K=1000\) & none & 3.44 & 12.72 \\
& layer 22, \(K=500\) & BPE 1000 & 1.85 & 9.78 \\
& layer 22, \(K=1000\) & BPE 2000 & 1.97 & 9.11 \\
& layer 22, \(K=2000\) & BPE 6000 & 1.47 & 7.68 \\
\midrule
\multicolumn{5}{@{}l}{\textit{B. Speech-unit span-corruption pretraining}} \\
& 0 epochs & -- & 3.61 & 14.189 \\
& 3 epochs & -- & 4.54 & 15.73 \\
& 10 epochs & -- & \textbf{5.81} & \textbf{18.452} \\
\bottomrule
\end{tabular}
\caption{Core design study; values are percentages at their reported
precision. The evaluation split is not recorded per ablation row, and the
10-epoch value equals the headline test result. The records also do not
identify which granularity runs omit speech-unit pretraining.}
\label{tab:core-ablation}
\end{table*}

\subsection{What works and what does not}
\label{sec:analysis}

\paragraph{Acoustic units and BPE.}
The reported unit rows ask whether later HuBERT representations and different
codebook sizes yield more useful inputs. At \(K=500\), the layer-22 row scores
above the layer-7 row: Hit@1 changes from 1.20 to 3.61 and Hit@20 from 6.00 to
14.189; among the uncompressed layer-22 rows, \(K=128\) has the best Hit@1
(3.76), whereas \(K=500\) has the best Hit@20 (14.189). Among these reported
configurations, later-layer units are associated with higher scores, while
codebook size trades top-rank and top-20 performance.

BPE was tested under the hypothesis that shorter acoustic-token sequences
would be easier for the encoder--decoder to process. In the matched \(K=500\)
comparison, BPE lowers Hit@1/Hit@20 from 3.61/14.189 to 1.85/9.78; at
\(K=1000\), it lowers them from 3.44/12.72 to 1.97/9.11. The tested BPE
settings therefore score below their same-\(K\) no-BPE rows, so the final
configuration omits BPE. The experiment does not isolate whether the
difference comes from information loss, optimization, or an unresolved shared
setting.

\paragraph{Speech-unit pretraining.}
Span-corruption pretraining tests whether unsupervised exposure can adapt the
text-pretrained backbone to acoustic units before retrieval fine-tuning. The
reported scores increase across 0, 3, and 10 epochs, from 3.61/14.189 to
4.54/15.73 and 5.81/18.452. The full 0-to-10 difference is \(+2.20\) Hit@1
and \(+4.263\) Hit@20. This is the strongest positive pattern in the report,
but the unresolved per-row split and shared settings prevent a causal or
test-set interpretation.

\begin{table}[t]
\centering
\small
\setlength{\tabcolsep}{5pt}
\begin{tabular}{@{}lllcc@{}}
\toprule
Study & Configuration & Intervention & Hit@1 & Hit@20 \\
\midrule
Identifier & string identifier & multi-token corpus ID & 3.61 & \textbf{14.19} \\
& atomic identifier & distinct unstructured ID & \textbf{4.54} & 13.69 \\
\midrule
Ranking loss & CE baseline & none & 3.61 & 14.189 \\
& CE + ranking & pairwise margin, in-batch negatives & \(\Delta{\approx}{-}0.88\) & \(\Delta{\approx}{-}1.349\) \\
\midrule
Compression & no Q-Former & none & 3.61 & 14.189 \\
& window Q-Former & \(L=17,\ Q=1\) & 2.73 & 12.584 \\
\bottomrule
\end{tabular}
\caption{Secondary design studies; values are percentages. Each intervention
is compared only with the baseline in the same reported subsection. The split
and pretraining setting are not recorded per row. The ranking-loss row gives
the reported approximate changes from its baseline, not reconstructed absolute
scores.}
\label{tab:secondary-ablation}
\end{table}

\paragraph{Identifier design.}
Atomic identifiers test whether eliminating multi-token decoding errors helps
retrieval. They increase Hit@1 from 3.61 to 4.54 but reduce Hit@20 from 14.19
to 13.69 relative to string identifiers. Identifier choice is therefore a
metric-dependent trade-off in this experiment, not a uniformly beneficial
change; the final system retains strings to favor the top-20 metric.

\paragraph{Ranking loss.}
A pairwise margin objective with in-batch negatives was intended to align
training more directly with ranking~\citep{Zeng2023ScalableAE,Sun2023LearningTT}.
Relative to the 3.61/14.189 cross-entropy baseline, the report gives
approximate decreases of 0.88 Hit@1 and 1.349 Hit@20, although convergence is
reported as faster. Under this negative-sampling and weighting setup, the
reported faster convergence does not accompany a higher final score, so the
final system keeps cross-entropy alone.

\paragraph{Windowed Q-Former.}
The windowed Q-Former tests whether compressing long acoustic-token sequences
before the T5 encoder improves retrieval~\citep{Pan2023COSMICDE,Tang2023SALMONNTG}.
With window length \(17\) and one learned query, it lowers Hit@1/Hit@20 from
3.61/14.189 to 2.73/12.584. This particular compression configuration scores
below its reported baseline; whether the cause is information loss or
insufficient optimization is not isolated by the experiment.
